\documentclass[10pt,a4paper]{amsart}
\usepackage{amsmath,amssymb,amsthm,graphicx}
\usepackage{verbatim}
\newenvironment{reference}{\comment}{\endcomment}
%\newenvironment{reference}{}{}
\newenvironment{slogan}{\comment}{\endcomment}
\newenvironment{history}{\comment}{\endcomment}
% For commutative diagrams we use Xy-pic
\usepackage[all]{xy}

% We use 2cell for 2-commutative diagrams.
\xyoption{2cell}
% Theorem environments.
%
\theoremstyle{plain}
\newtheorem{theorem}[subsection]{Theorem}
\newtheorem{proposition}[subsection]{Proposition}
\newtheorem{lemma}[subsection]{Lemma}

\theoremstyle{definition}
\newtheorem{definition}[subsection]{Definition}
\newtheorem{example}[subsection]{Example}
\newtheorem{exercise}[subsection]{Exercise}
\newtheorem{situation}[subsection]{Situation}

\theoremstyle{remark}
\newtheorem{remark}[subsection]{Remark}
\newtheorem{remarks}[subsection]{Remarks}

\numberwithin{equation}{subsection}

% Macros
%
\def\lim{\mathop{\mathrm{lim}}\nolimits}
\def\colim{\mathop{\mathrm{colim}}\nolimits}
\def\Spec{\mathop{\mathrm{Spec}}}
\def\Hom{\mathop{\mathrm{Hom}}\nolimits}
\def\Ext{\mathop{\mathrm{Ext}}\nolimits}
\def\SheafHom{\mathop{\mathcal{H}\!\mathit{om}}\nolimits}
\def\SheafExt{\mathop{\mathcal{E}\!\mathit{xt}}\nolimits}
\def\Sch{\mathit{Sch}}
\def\Mor{\mathop{\mathrm{Mor}}\nolimits}
\def\Ob{\mathop{\mathrm{Ob}}\nolimits}
\def\Sh{\mathop{\mathit{Sh}}\nolimits}
\def\NL{\mathop{N\!L}\nolimits}
\def\CH{\mathop{\mathrm{CH}}\nolimits}
\def\proetale{{pro\text{-}\acute{e}tale}}
\def\etale{{\acute{e}tale}}
\def\QCoh{\mathit{QCoh}}
\def\Ker{\mathop{\mathrm{Ker}}}
\def\Im{\mathop{\mathrm{Im}}}
\def\Coker{\mathop{\mathrm{Coker}}}
\def\Coim{\mathop{\mathrm{Coim}}}

% Boxtimes
%
\DeclareMathSymbol{\boxtimes}{\mathbin}{AMSa}{"02}

%
% Macros for moduli stacks/spaces
%
\def\QCohstack{\mathcal{QC}\!\mathit{oh}}
\def\Cohstack{\mathcal{C}\!\mathit{oh}}
\def\Spacesstack{\mathcal{S}\!\mathit{paces}}
\def\Quotfunctor{\mathrm{Quot}}
\def\Hilbfunctor{\mathrm{Hilb}}
\def\Curvesstack{\mathcal{C}\!\mathit{urves}}
\def\Polarizedstack{\mathcal{P}\!\mathit{olarized}}
\def\Complexesstack{\mathcal{C}\!\mathit{omplexes}}
% \Pic is the operator that assigns to X its picard group, usage \Pic(X)
% \Picardstack_{X/B} denotes the Picard stack of X over B
% \Picardfunctor_{X/B} denotes the Picard functor of X over B
\def\Pic{\mathop{\mathrm{Pic}}\nolimits}
\def\Picardstack{\mathcal{P}\!\mathit{ic}}
\def\Picardfunctor{\mathrm{Pic}}
\def\Deformationcategory{\mathcal{D}\!\mathit{ef}}
\usepackage{cleveref}
\crefname{originallemma}{original lemma}{original lemmas}
\Crefname{originallemma}{Original lemma}{Original lemmas}
\makeatletter
\newlabel{lemma-local-CM}{{lemma-local-CM}{source}{}{sourcepointer-1}{}}
\newlabel{lemma-local-CM@cref}{{[originallemma][0][]lemma-local-CM}{[1][0][]source}{}{}{}}
\makeatother
\title{The locus where a module is Cohen-Macaulay is open and its closure is the support}
\author{The Stacks Project Authors}
\date{stacks-project:dualizing.tex:local-snapshot-20261012; GNU FDL 1.2 or later}
\begin{document}
\maketitle
\noindent\textit{Editorial note.} All delivered mathematical text is the author's own, copied from the pinned Stacks Project chapter source sources/stacks-dualizing/pinned-original-source.tex; no statement or proof step was rewritten, reordered or removed.\par\medskip
\noindent\textit{Editorial note.} Preamble provenance: the Stacks Project keeps one single shared preamble for every chapter; this card takes it directly from the staged sources/stacks-dualizing/preamble.tex as three verbatim spans (lines 8-12, 14-18 and 154-217), so no derived file was written and nothing was added or redefined. The wrapper class and the base packages are supplied by the assembler.\par\medskip
\noindent\textit{Editorial note.} Bibliography provenance: neither the counted statement nor its proof contains a citation (this lane rescans the whole span before writing the card), so the card carries no bibliography and no entry was reconstructed.\par\medskip
\noindent\textit{Editorial note.} Selection provenance: this card was opened under the file-based control contract (CONTROL.md), with the mechanical pre-card gate run on its statement and proof interval before the card was written.\par\medskip
\noindent\textit{Allowed standard prerequisites (editorial).}\begin{itemize}
\item The statement that a module over a local ring with vanishing local cohomology of the dual in all but one degree is Cohen-Macaulay (and conversely). Invoked inside the counted proof (source lines 3893-3898). Identifies the locus with the set of points where the dual has a unique nonzero cohomology module.
\item Finiteness of the cohomology modules of the dual and the vanishing criterion for finitely generated modules. Used inside the counted proof (source lines 3899-3902). Produces a localizing element killing the unwanted cohomology.
\end{itemize}\medskip
\ifdefined\hypertarget\hypertarget{sourcepointer-1}{}\fi
\noindent\textit{Source reference.} Original lemma lemma-local-CM: sources/stacks-dualizing/pinned-original-source.tex, source lines 3307--3307. Navigation pointer only: the counted text refers to the Stacks item lemma-local-CM, which is an existing labeled result of this same chapter and is not part of this card. The pointer carries the target's own label as its number and its definition line as its span; its content is not reproduced here.\par\medskip
\begin{lemma}
\label{lemma-CM-open}
Let $A$ be a Noetherian ring with dualizing complex $\omega_A^\bullet$.
Let $M$ be a finite $A$-module. Then
$$
U = \{\mathfrak p \in \Spec(A) \mid M_\mathfrak p\text{ is Cohen-Macaulay}\}
$$
is an open subset of $\Spec(A)$ whose intersection with
$\text{Supp}(M)$ is dense.
\end{lemma}

\begin{proof}
If $\mathfrak p$ is a generic point of $\text{Supp}(M)$, then
$\text{depth}(M_\mathfrak p) = \dim(M_\mathfrak p) = 0$
and hence $\mathfrak p \in U$. This proves denseness.
If $\mathfrak p \in U$, then we see that
$$
R\Hom_A(M, \omega_A^\bullet)_\mathfrak p =
R\Hom_{A_\mathfrak p}(M_\mathfrak p, (\omega_A^\bullet)_\mathfrak p)
$$
has a unique nonzero cohomology module, say in degree $i_0$, by
Lemma \ref{lemma-local-CM}.
Since $R\Hom_A(M, \omega_A^\bullet)$
has only a finite number of nonzero cohomology modules $H^i$
and since each of these is a finite $A$-module, we can
find an $f \in A$, $f \not \in \mathfrak p$ such that
$(H^i)_f = 0$ for $i \not = i_0$. Then
$R\Hom_A(M, \omega_A^\bullet)_f$ has a unique nonzero cohomology
module and reversing the arguments just given we find
that $D(f) \subset U$.
\end{proof}

\end{document}
